% !TEX root = ../monograph.tex
\section{From the current identity to coefficient estimates}
\label{ch17:marked-rates}

The integration by parts needed here is performed before the
large-field expansion. This order matters: a large-field activity is
only measurable, so its derivative cannot be inserted into the
expansion as an ordinary bounded function.

\subsection{A differentiated block kernel}

In the exact kernel $k^{(L)}_\beta$ of \eqref{ch04:kernel}, condition
on the fine field $U$ and all auxiliary choices $\alpha$. The conditional density of $V_e$ is
$q_{b\beta}(V_e Q_e^{-1})$, where
$q_\gamma(u)=Z_\gamma^{-1}\exp[-\gamma w(u)]$.
For every fixed order $r$, compactness and smoothness of $w$ give
\begin{equation}\label{ch17:kernel-score}
 |D^r q_\gamma(u)|\le C_r(1+\gamma)^r q_\gamma(u).
\end{equation}
Here $D^r$ denotes $r$ successive left derivatives in unit Lie algebra
directions. Dividing by $q_\gamma$, the result is a polynomial in
derivatives of $w$ with at most $r$ factors of $\gamma$. The normalization is
independent of $u$. For a product of link kernels the corresponding
bound has an additional polynomial in the number of differentiated
links. No lower bound on $q_\gamma$ is used.

First truncate the exponentially local current and the kernels of the critical-configuration map (the coefficients of its Taylor expansion in the coarse field) on a finite neighborhood. For the vector field use a finite
gauge-covariant extension of its controlled jets through field order four,
not the uncontrolled higher derivatives of a deep transferred current.
Use the coefficientwise gauge-covariant extension: replace each
fixed-order curvature or covariant-difference monomial by its plaquette
and contour-variable representative, and correct its lower jets
triangularly through order four. Each fixed number of link derivatives
adds only a fixed power of the contour length. The extension therefore
has a bounded weighted norm at a strictly narrower exponential width;
it is not an arbitrary jet section on a space of growing dimension.
Complete in that weighted norm first and then truncate. All additional
fixed-order derivative bounds depend only on the controlled jets and
the fixed stage. Extend this local small-field formula smoothly to
the compact product group. One can use a smooth cutoff
supported strictly inside the logarithmic chart, equal to one on the
smaller chart used for Taylor extraction, and then average over the
compact local gauge group. Equivalently extend its finite covariant
holonomy jets. This preserves the required jets and covariance.
The finite derivative factors are inverse powers of the chart width and
polynomials in the support size. They are allowed in the bounds below.
Only after this extension do we integrate by parts. For an exponentially
local current do this first for finite support truncations. Fixed powers
of the support length are absorbed by a strict margin in the
exponential coefficient width. The derivative bounds and tail errors
are consequently summable, so the finite-period identity passes to
the weighted tail limit.

Let $k(U,V)$ be that exact kernel, let $f(U)$ be the entire incoming
density, and put $F(V)=\int k(U,V)f(U)\,dU$. All original fine large-field indicators
and compatibility conditions are part of $f$ or of the integral over the local gauge transformations;
none is differentiated. For a smooth coarse vector field $Y$, define
\[
 \mathcal D_Y k=(\operatorname{div}_{\rm Haar}Y)k+D_Vk[Y].
\]
Dominated differentiation, using \eqref{ch17:kernel-score}, gives the
\emph{density identity}
\begin{equation}\label{ch17:kernel-divergence}
 \operatorname{div}_{\rm Haar}(FY)
       =\int\mathcal D_Y k(U,V)f(U)\,dU.
\end{equation}
Consequently its integral against every spectator independent of the
varied coarse links is zero. This is an identity for the full differentiated kernel,
including its nonpolynomial remainder, not an assertion that a finite
Taylor polynomial alone has zero exterior moments, that is, zero integral against every spectator independent of the varied links.

The right hand side is an admissible conditional integral with the differentiated kernel inserted.
Its extra factor is a finite sum of smooth local kernel scores times
$Y$ or its derivative. On a protected dependence set the incoming
measurable product remains fixed under the change of variables in
Chapter~\ref{ch06:large-fields}. Additional coarse derivatives act
on the score and the common density ratio only. The bounded pointwise-jet estimate
\eqref{ch06:actual-jet-integral} therefore applies with a fixed
polynomial factor in $\beta$, support size and current derivative bounds.
Rooting the connected expansion at this extra factor gives the same
expansion with insertions as in Chapter~\ref{ch11:insertions}. In particular
large-field terms retain their original prescribed large-field sets.
This argument does not differentiate the event indicators introduced when that
conditional integral is subsequently decomposed.

The small-field contribution needs a sharper identity than the absolute
score bound. Fix the full dependence set of a rooted connected component and use one
additional coarse derivative beyond those prescribed for the extracted
coefficient. On its smaller small-field chart, make exactly the conditional
Gaussian change of variables used to construct the small-field expansion.
Write its normalized conditional integrand as $p(V,z)$ and its integral
as $Z_{\rm ord}(V)$. For finite small-field selections, differentiation
under this integral gives
\[
 D_VZ_{\rm ord}[Y]=\int D_Vp[Y]\,dz,
 \qquad
 \int D_Vp[Y]\,dz
 =Z_{\rm ord}D_V\log Z_{\rm ord}[Y].
\]
Every Jacobian, relative determinant and conditional normalization is
included in $p$; thus this is not the derivative of an unnormalized
Gaussian moment. The absolutely convergent small-field logarithm, with
one additional derivative, passes this identity to the countable expansion.
By its definition in Chapter~\ref{ch09:renormalization},\mgubi{Issue 9: $Z_{\rm ord}$ is defined only above, for one component; $E_{\rm reg}$ is defined nowhere. Chapter~\ref{ch09:renormalization} gives this form with a constant only for related objects (the pure small-field block integral; the constant extracted once, globally; part (a) of Assumption~\ref{ch09:forest-assumption}, unproved). A constant for one component needs an argument.}
\[
 Z_{\rm ord}=\mathcal C
       \exp[-\beta S-\mathcal L-E_{\rm reg}],
 \qquad D_V\log\mathcal C=0.
\]
If an intermediate local normalization depends on $V$, it is included
in $\mathcal L+E_{\rm reg}$, not in $\mathcal C$. Consequently the
small-field score is \emph{exactly}
$-\beta DS[Y]-D\mathcal L[Y]-DE_{\rm reg}[Y]$.
The regular term is estimated using its already proved differentiated
regular norm. We do not estimate it by the polynomial-in-$\beta$
absolute score bound.

This argument uses the full rooted component, not just the score
stencil. In the rooted connected expansion, small-field factors in
components disjoint from the component of the differentiated kernel cancel
between the differentiated numerator and its normalization. Every term with a large fluctuation-field event anywhere in the connected
component of the differentiated kernel, including a distant such event joined to the score, is kept in the complementary
sum. The smooth cutoffs used above are constant on the smaller small-field
chart; their derivatives are supported where at least one enlarged
chart condition fails. Attach this condition to its dependence set. The same
Gaussian or energy estimate for the nested cutoffs bounds this new
large fluctuation-field event, while all incoming prescribed large-field sets are retained.
The absolute score bound is used only for these complementary terms,
where the exponentially small probability of the large fluctuation-field event absorbs its polynomial
$\beta$ factor. No derivative of a hard indicator is taken. This proves
the small-field cancellation and the complementary large-field estimate
separately, which is essential to preserve the regular power gain.

\subsection{The finite degree-five coefficient}

At the comparison scale $k$ write
\[
 h=(\log\log\widehat\beta_k)^2,\qquad H=e^h,
 \qquad t=\widehat\beta_k^{-1/2}H.
\]
All constants in this paragraph are for a fixed stage and the fixed
finite derivative orders and degrees; they are fixed before its lower bound on the
reference coupling. Let $\mathcal P_5$ denote radial Taylor extraction through
degree five and let $f$ be the number of curvature factors of a homogeneous
coefficient. Set $\mathcal P_{5,t}=t^{f-2}\mathcal P_5$, followed by
real signed-permutation averaging. There are no invariant linear forms on $\gLie$. Thus the nonzero vector coefficients under consideration
have $f\ge2$, and this normalization does not introduce a negative
power of $t$ beyond the displayed $t^{-2}$.

Apply the exact step with insertions to \eqref{ch17:kernel-divergence}, retaining
its entire remainder, and use the current identity through order four.
Let $j_k$ be the radial coefficient of the current divergence, in the
same threshold coordinates. The classical term $-\beta DS[Y]$ is
$-\beta j_k$ up to the contracting current defect. Thus the coefficient
of the full differentiated kernel, with zero exterior moments, is $\Sigma_k-\beta j_k$;
$\Sigma_k$ alone is not asserted to have zero exterior moments. Its
remaining degree-five scalar-vector coefficient can be written
\begin{equation}\label{ch17:finite-sigma}
 \Sigma_k=\mathcal P_{5,t}
       (\operatorname{div}_{\rm Haar}Y-D\mathcal L[Y])
       +E_{\rm stat}+E_{\rm tail}+E_{\rm reg}+E_{\rm mark}.
\end{equation}
Here the kernels of the critical-configuration map and the loop kernels are truncated at
$C_j(\log\widehat\beta_k)^2$. Their common exponential decay makes
$E_{\rm stat}$ and $E_{\rm tail}$ at most $H^{-1}$ after finite
projection factors. The regular error satisfies
\[
 t^{-2}\epsilon_k^{\rm reg}
       =\widehat\beta_k^{-1/10}H^{-2}.
\]
A directional derivative along $Y$ has the same bound up to the fixed
current norm, since the derivative bound contributes a factor $t^{-1}$
and $Y=O(t)$ on the small-field chart. The previously proved logarithmic current bounds
and fixed localization factors are absorbed by increasing the lower bound on the reference coupling.

For the remainder with insertions one must first cancel the whole classical
current divergence, not merely use the absolute kernel-score bound.
The current identity agrees through field order four. The finite
gauge-covariant extension just chosen has controlled higher derivatives,
so the remaining classical error starts at order five and contributes
\[
 t^{-2}\beta t^5=O(\widehat\beta_k^{-1/2}H^3).
\]
Fixed current and localization factors preserve a bound of the form
$C_j\widehat\beta_k^{-1/10}H^{C_j}$ after increasing the lower bound on the reference coupling. The small-field loop
and Jacobian Taylor remainders have no leading factor $\beta$ and
gain an additional small-field factor. This controls the high-field-order remainder. It does not control a
low-field-order term with many spatial differences: vanishing Taylor coefficients through the specified scaling dimension alone give no additional small-field factor. The remaining low-field-order spatial terms are controlled by exact
polynomial preservation, rather than by an extra power of the field.

\begin{proposition}[Low-moment estimate for spatial remainders]
\label{ch17:spatial-moments}
For the linearization of the critical configuration and the regular local terms of the differentiated kernel just constructed,
a remainder whose radial Taylor coefficients through scaling dimension five vanish
has extracted degree-five coefficient $O(H^{-1/2})$. The estimate is
uniform in period and in the finite current truncations, after the
stage orders and thresholds have been chosen with one extra derivative
order. It applies to both the nonclassical coefficient and the
normalized current divergence.
\end{proposition}
\begin{proof}
First identify the linearization of the critical configuration on polynomials. Recall the notation of
Chapter~\ref{ch05:gaussian}: $D_l,D$ are the rescaled fine and coarse
covariance symbols, $x_l,x$ their gradient columns, and
$\alpha=\sum_l|m_l|^2\tr D_l$.
Write $s=\lambda_{M_c}(p)$ and
$D=s^{-1}\operatorname{diag}(B_a)$, where the $B_a$ are the
regularized diagonal covariance factors. The scalar correction
\[
 \chi=\frac{x^*\operatorname{diag}(B_a^{-1})x}{s}
\]
is analytic and nonzero at zero, as in the pole cancellation for the
formulas for the linearization. Put
$\sigma=\chi^{-1}x^*\operatorname{diag}(B_a^{-1})$, so that
$\pi=s^{-1}\sigma$ in those formulas. For a nonzero alias define
\[
 g_{l,a}=\overline{m_l}
 \left(\frac{\tr D_l}{s\alpha}
       -\frac{(D_l)_{aa}|r_L(\theta_{l,a})|^2}{B_a}\right).
\]
Substitution into the linearization built with the linear averaging operator gives the regular formula
\[
 \mathcal A_l=\mathcal B_l+
          \operatorname{diag}(g_{l,a})x_l\sigma,
 \qquad \mathcal B_l=D_lF_{1,l}^*D^{-1}.
\]
Here $D_l$ is regular at a nonzero alias, and $s\alpha$ and
$B_a$ are analytic and nonzero at zero. Thus no uncancelled factor
$s^{-1}$ remains. Both terms contain $\overline{m_l}$, with
$m_l=w_L(\theta_l)^j$ and $w_L(\theta_l(0))=0$.
Every derivative of $\mathcal A_l$ of order less than $j$ therefore
vanishes at zero. The column at the zero alias has value $I$. Through \eqref{ch17:polynomial-response}, this
order-$j$ vanishing is also what the preservation of the order of vanishing of the
current defect in Section~\ref{ch17:quadratic-current} uses. Stars and conjugates in these symbols mean analytic
transpose at the opposite momentum, as in Chapter~\ref{ch05:gaussian}.

Let $A_x^{ab}(z)$ be the physical row from coarse direction $b$ at
$\lfloor x/L\rfloor+z$ to fine direction $a$ at $x$.
Differentiating the Fourier inversion formula through order less than $j$
therefore gives, for every vector polynomial $p$ of degree less than $j$,
\begin{equation}\label{ch17:polynomial-response}
 (W_sp)(x)_a=\sum_{b,z}A_0^{ab}(z)p_b(x/L+z),
 \qquad \sum_z A_0^{ab}(z)=L^{-1}\delta_{ab}.
\end{equation}
The exponential row bound justifies the derivatives and the sum.
Polarization gives mixed polynomial moments. In particular the output
is a global polynomial of no larger degree, not a polynomial separately
on each block. Its homogeneous degree-$r$ leading term is
$L^{-r-1}p(x)$; lower coefficients obey the same ratio-independent row bounds.

Choose $j$ above the order of the coefficient evaluations. In a polarized coefficient
evaluation take only tuples of potential polynomials whose summed scaling
dimension is at most five. Equation~\eqref{ch17:polynomial-response}
preserves or lowers each polynomial degree, so every resulting tuple
remains within that same finite list. The polynomial-moment
cancellation of Section~\ref{ch11:one-graph} makes its deterministic
coefficient evaluation exactly zero. This is not a claim of vanishing on arbitrary
polynomial inputs of degree less than $j$. Truncating the coarse field to the ball of radius
$C_j(\log\widehat\beta_k)^2$ changes this coefficient evaluation only by the
exponential row tail, times fixed polynomial directional and support
factors. That tail is smaller than every prescribed power of $H^{-1}$.

The smooth representative of the linearization of the critical configuration is
$W_a=W_s-d_fERW_s$. Its difference from $W_s$ is a fine gradient.
The coefficient of these regular terms is a spacetime vector with trivial adjoint action on its values,
invariant under local fine gauge transformations. Indeed gauge
equivariance of $Y$ and invariance of Haar volume imply invariance of
$\operatorname{div}_{\rm Haar}Y$; the action and loop variations
are invariant for the same reason. Apply the corresponding fine gauge
transformation before evaluating its local jet. The Baker--Campbell--Hausdorff formula with the same local gauge choice compares the two exponential
representatives with an error $C_jP(B)t^2$ in the link logarithms,
where $P$ is a fixed polynomial. Dividing by the fine threshold gives
$C_jP(B)t$. This argument uses the exact correction to the averaging operation, not an
unmatched linear averaging operator.

It remains to compare the conditional insertion itself to this
polynomial coefficient evaluation. Normalize the input link logarithms by $t$.
The nonlinear critical configuration, the same local gauge coordinate change,
and the deterministic truncation contribute $C_jP(B)t$ and the exponential
tail just described. Gaussian fluctuations contribute
$C_jP(B)g/t$, where $g=\beta^{-1/2}$, by the fundamental theorem of
calculus and the one additional derivative. Explicitly, if
$p(z)\ge1$ is a polynomial majorant for the finite Gaussian coordinates,
the difference between the conditional and the deterministic evaluation in
each prescribed derivative is bounded by
\[
 C_j N_k P(B)\bigl(t+g/t+e^{-c_jB}\bigr)\,p(z)^{C_j},
\]
with an appropriate enlargement of the truncation radius in the last term.
Here $N_k$ is the norm of the already $t^{-2}$-normalized regular
term of the differentiated kernel and its derivatives in normalized link logarithms, including
its current and loop coefficients. The extraction normalized by the number of curvature factors has
already been performed; no further $t^{-2}$ is applied to this bound. For this bound subtract the deterministic coefficient evaluation
inside the same normalized conditional integral. Its derivative
through the prescribed scaling dimension is zero. Use the joint
polynomial moments and differentiated density scores of that one
conditional law to estimate the evaluation difference; no small
total-variation distance between the nonlinear law and a Gaussian
is asserted. The pointwise jet and moment estimates are those of \eqref{ch06:actual-jet-integral};
no differentiation of incoming large-field factors is involved.
The complement is estimated by the already retained large-field bound. Finite coefficient recovery adds a factor depending only on the stage.

There are no invariant terms that are constant spacetime vectors or linear in the $\gLie$-valued field. Hence these regular terms of the differentiated kernel start at field order two. Their
values and normalized-link derivatives on the threshold ball contain
a factor $t^2$; the explicit $t^{-2}$ removes exactly this factor.
The classical part has already been cancelled and its order-five
remainder treated separately above. Thus for the present regular terms
$N_k\le H^{\theta_j+o(1)}$, with
$\theta_j\to0$: current norms are logarithmic and the retained loop
norm has this bound. Since $B$ and all truncation radii are fixed powers
of $\log\widehat\beta_k$, their polynomial factors are $H^{o(1)}$.
The coupling band gives $g/t\le C H^{-1}$, whereas $t$ and
the truncation tail decay faster than every fixed inverse power of $H$.
The resulting coefficient error is $H^{-1+\theta_j+o(1)}$, hence
at most $CH^{-1/2}$ after increasing the lower bound on the reference coupling. Large supports are first truncated
using a margin in the support exponent; its tail obeys the same
bound. This proves the proposition for the full local sum.
\end{proof}

The differentiated-kernel
bound \eqref{ch17:kernel-score} applies to the complementary part.
The complementary part retains the square root
$\exp(-H^{1/5}/2)$ of the bound \eqref{ch04:activity-bounds} on the large-field activities. Combining Proposition~\ref{ch17:spatial-moments} with the high-field
remainder, the nonwrapping contributions are bounded by
\begin{equation}\label{ch17:sigma-error}
 C H^{-1/2}+C_j\widehat\beta_k^{-1/10}H^{C_j}
       +C_j\widehat\beta_k^{C_j}H^{C_j}e^{-H^{1/5}/2}.
\end{equation}
The fixed-order score has changed the polynomial constants, not the
eligibility event \eqref{ch04:eligibility} of a large-field activity or its exponential gain. If $N_{\rm per}$ is the
period in scale-$k$ units, a wrapping term has covering number at least
$N_{\rm per}/[10(2B_k+1)]-1$. Retaining two units of the support exponent gives the
additional error
\[
 C_j\widehat\beta_k^{C_j}H^{C_j}
       \exp\bigl(-N_{\rm per}/[5(2B_k+1)]+2\bigr).
\]
It vanishes in the ordered period limit at fixed $k$.

The two terms other than $CH^{-1/2}$ in
\eqref{ch17:sigma-error} are $o(H^{-1})$. To check the
order of limits explicitly, put $x=\log\log\widehat\beta_k$.
The logarithm of the first bound multiplied by $H$ is at most
$C_j+(C_j+1)x^2-e^x/10$; that of the second is at most
$C_j+C_je^x+(C_j+1)x^2-e^{x^2/5}/2$.
Both tend to minus infinity at each fixed stage. Increasing its lower bound on the
reference coupling makes the bounds uniform on that stage.

The loop coefficient has zero $\rho$ by its two moment cancellations:
its quadratic kernel annihilates both constant and affine links, whereas
the cubic momentum coefficient functional pairs a constant with a quadratic link or two
affine links. This uses the loop kernel of Chapter~\ref{ch07:matching},
not only gauge invariance.
The calculation with mass regularization identifies the remaining coefficient as
$\mathfrak h(Y)/12$. One can check the factor directly in the three
vector coordinates of \eqref{ch17:rho}, before the quotient by the equation of motion: the cubic momentum
functional takes values $-24,-96,0$, while $\rho$ takes values
$-2,-8,0$. The lattice correction is the matching lower-order
Taylor correction in $\mathfrak h$.

Choose the number of transfers and then the period at each fixed $k$ so
that the error of the mass regularization and the finite momentum approximation are
smaller than $\widehat\beta_k^{-30}$. The retained loop bound is
$2H^{\theta_j}$, with $\theta_j\to0$, and current derivative norms
are logarithmic. Sufficiently large lower bounds on the reference coupling therefore give joint ordered limit
coefficients satisfying
\begin{equation}\label{ch17:sigma-bound}
 \|\Sigma_k\|\le C H^{1/6},\qquad
 \rho(\Sigma_k)\longrightarrow-\kappa_G.
\end{equation}
All coefficient maps and the underlying correlations are retained
in the same joint limit. Compactness is used only after the uniform
finite-cutoff estimates have been established.

\subsection{The kernel vector and the necessary rates}

Let $j_k$ be the degree-five coefficient of the transferred current's
divergence, in the same threshold coordinates. Its local Taylor
coefficient has zero $\rho$: the trace is a forward difference of a
degree-three primitive, whose invariant projection is a multiple of the identity field.
The same localization and error bounds with insertions compare the transferred
coefficient with this Taylor coefficient within $CH^{-1/4}$.
Consequently, in the notation of the stability argument, with $v=(0,0,1)$ the kernel vector of the free quotient $s$,
\[
 j_i\to-v,\quad |\rho(j_i)|\le CH_i^{-1/4},\qquad
 v_i^{\rm eq}=\beta_i^{-1}\Sigma_{k_i}-j_i\to v.
\]
With $e_i=H_i^{1/5}$ this gives $e_i\rho(v_i^{\rm eq})\to0$.

There is also a stronger, independent bound on the image of this
vector. The differentiated kernel is carried as a density-valued factor, without
dividing it by $F$ or claiming that the resulting quotient is local.
More directly, for every incoming coefficient $f_A(U)$ with insertions,
including all the original fine spectator factors,
\[
 \int \mathcal D_Y k(U,V)\,dV=0
 \quad\hbox{for each }U,\qquad
 \int \mathcal D_Y k(U,V)f_A(U)\,dU\,dV=0.
\]
The same differentiated kernel is used for every $A$, so the joint rooted
expansion retains all pushed spectator cross-modes and this exact
zero-moment statement. There is no smoothness requirement on $f_A$.
For a smooth local coarse spectator $A(V)$ one also has
\[
 \int A(V)\!\int\mathcal D_Y k(U,V)f(U)\,dU\,dV
 =-\int D_VA[Y]\!\int k(U,V)f(U)\,dU\,dV.
\]
It is zero when the spectator does not read the varied links. For an
exponentially local spectator that need not have finite range, first truncate its dependence set
and use this identity on disjoint terms. Retain the remaining terms
in the same joint expansion with insertions. A connected term reaching from
the differentiated kernel to a spectator at distance $R$ has covering number at least
$R/(2B_k+1)-1$. Retaining two units of the support exponent bounds these terms by
$C\exp[-2R/(2B_k+1)]$, with their normalized large-field majorants as well.
For spectator configurations separated by a fixed positive multiple of $a_n$, their
distance in scale-$k(n)$ units is at least $c\exp(T_n)$. This tail
is therefore smaller than $e^{-T_n/100}$ after all fixed factors.
Every spectator subset uses the same density without insertions.
Write $\tau_j=\log(1/a_j)$ and
$T_n=(\log\widehat\beta_n)^{48}$. Choose $k(n)$ to be the smallest
integer $k\ge n$ for which $\tau_k-\tau_n\ge T_n$. Match the full
coefficient and remainder backwards from $k(n)$ to $n$. Equation~\eqref{ch11:descent-error} gives
$e^{-T_n/8}$ before multiplication by its constants. Fixed projection and
threshold factors have logarithm $O(\log\widehat\beta_n)$; accumulated
ratio-independent coefficient-map factors are $o(T_n)$. Here it is only the latter
constants that are multiplied along the interval of scales. Ratio-dependent nonlinear
jet constants have already been absorbed into the errors controlled by the lower bounds on the reference coupling of the stages.

Place a copy in each reflected cell at the first backward crossing of $B_n^2$, so
$B_n^2\le a_m/a_n\le L_*B_n^2$. The lower bound on the partition function
in \eqref{ch11:replicated-bound} contributes at most
\[
 L_*^4B_n^8q_n+6=O((\log\widehat\beta_n)^{22})=o(T_n),
\]
using $L_*\le\log\widehat\beta_n$, $B_n=O((\log\widehat\beta_n)^2)$
and the per-site normalization exponent $q_n=O((\log\widehat\beta_n)^2)$.
The inverse source weights in \eqref{ch11:replicated-bound} have logarithm $O(\log\widehat\beta_n)$.
The dimensionful normalization of the insertion cancels the factor $a_n^{-5}$.
Thus the normalized residual of \eqref{ch11:descent-error} is at most $e^{-T_n/100}$.
This argument uses the full differentiated kernel with zero exterior moments; its Taylor remainder
is carried by the graph of the backward comparison, not discarded.

In the choice of radii along the diagonal subsequence require
$r_i/a_{n_i}\le\exp((\log\log\widehat\beta_{n_i})^3)$.
Finite separating matrices of free two-point correlation pairings have uniformly bounded left
inverses. For every fixed $p$ the image therefore satisfies
\[
 b_i^p\|s_i(v_i^{\rm eq})\|
 \le C b_i^p\exp\bigl(5(\log\log\widehat\beta_{n_i})^3-T_{n_i}/100\bigr)
 \longrightarrow0.
\]
This proves the quantitative kernel hypothesis in the stability lemma.
Its bounded right inverses give bounded $\lambda_i$, with
$e_i\|\lambda_i s_i-\rho\|\to0$.

The same identity for the differentiated kernel gives
$\beta_iK_iB_i^L-s_i\Sigma_{k_i}\to0$.
For $Z=X,D$, the ordinary backward comparison gives
$K_iB_i^Z-s_i\operatorname{div}_{\rm raw}Z=o(b_i^{-p})$ for every
fixed $p$. Apply the bounded functional and use
$H_i^{1/6}/e_i\to0$ in \eqref{ch17:sigma-bound}.
This proves the first two lines of \eqref{ch17:transfer}.

Finally the trace of $X$ is identically zero, while the transferred
trace of $L$ is an exact forward difference. Match its degree-three
primitive at a fixed late pair of cell scale and stopping scale. The backward comparison gives endpoint
error $C e^{-(\tau_k-\tau_{k_0})/8}$, uniformly near the configurations used in the finite separating
correlation pairings. Subexponential coefficient conversions and support truncation
make this $O(a_k^\eta)$ for any fixed $0<\eta<1/8$.
Average over $m=\lfloor a_k^{-1+\eta/2}\rfloor$ consecutive differences.
The identity-field parts of the two endpoint values are equal and cancel, leaving at most
$2Ca_k^\eta/(ma_k)=O(a_k^{\eta/2})$.
The averaging length $ma_k$ in continuum units tends to zero, so the finite matrix of
correlation pairings converges to its original invertible matrix. This proves the
last rate in \eqref{ch17:transfer}, including multiplication by every
fixed power of $b_i$. Ordinary Taylor matching gives its remaining
convergence $A_iB_i^Z\to Z$.
